\documentclass[11pt,a4paper]{article}
\usepackage{import}
\subimport{../}{preamble.tex} 
\ungradedsolutionstrue
\gradedsolutionstrue
\chead{\textbf{Do-calculus and Faithfulness}}

\begin{document}



\section{Chain-Rule}

In this exercise, you will prove a well-known chain-rule for the case of Markov kernels. 
Let $K: \mathcal{X} \dashrightarrow \mathcal{Y} \times \mathcal{W}$ be any Markov kernel. 
Prove that 
\begin{equation*}
  K(Y | X) = K(Y | W, X) \circ K(W | X),
\end{equation*}
where $K(Y | X)$ and $K(W | X)$ are the marginal Markov kernels and $K(Y | W, X)$ is the conditional Markov kernel.
You should base your proof on the following definitions (and your general understanding), and make clear where you use which definition:
\begin{enumerate}
  \item By definition, for all measurable $D \subseteq \mathcal{Y}$ and $x \in \mathcal{X}$: 
    \begin{equation*}
      \left[ K(Y | W, X) \circ K(W | X)\right](D | x) = \int_{\mathcal{W}} K(Y \in D | W = w, X = x) K(W \in \text{d} w | X = x).
    \end{equation*}
  \item By definition, for all measurable $F \subseteq \mathcal{Y} \times \mathcal{W}$ and $x \in X$:
    \begin{equation*}
      \left[ K(Y | W, X) \otimes K(W | X)\right](F | x) = \int_{\mathcal{W}} \bigg[ \int_{\mathcal{Y}} \mathds{1}_{F}(y, w) K(Y \in \text{d} y | W = w, X = x) \bigg] K(W \in \text{d} w | X = x).
  \end{equation*}
\item  $K(Y \in D | X = x) \coloneq K(Y \in D, W \in \mathcal{W} | X = x)$, and similarly for $K(W \in E | X = x)$, for all measurable $D \subseteq \mathcal{Y}$ and $E \subseteq \mathcal{W}$ and all $x \in \mathcal{X}$.
\item The definition of conditional Markov kernels.  
  \end{enumerate} 

  \begin{gradedsolution}
    For all measurable $D \subseteq \mathcal{Y}$ and $x \in X$ we have:
\begin{align*}
  \big[ K(Y | X) \big] (D | x) & \overset{3.}{=} \big[ K(Y, W | X)\big] (D \times \mathcal{W} | x) \\
  & \overset{4.}{=} \big[ K(Y | W, X) \otimes K(W | X)\big] (D \times \mathcal{W} | x) \\
  & \overset{2.}{=} \int_{\mathcal{W}} \bigg[ \int_{\mathcal{Y}} \underbrace{\mathds{1}_{D \times \mathcal{W}}(y, w)}_{= \mathds{1}_D(y)} K(Y \in \text{d} y | W = w, X = x)\bigg] K(W \in \text{d} w | X = x) \\
  & = \int_{\mathcal{W}} K(Y \in D | W = w, X = x) K(W \in \text{d} w | X = x) \\
  & \overset{1.}{=} \big[  K(Y | W, X) \circ K(W | X)\big] (D | x).
\end{align*}
That proves the desired equality.
  \end{gradedsolution}



\section{Do-Calculus I}
\begin{figure}[ht]
  \centering
  \begin{tikzpicture}[scale=1, transform shape]
      \tikzstyle{every node} = [draw,shape=circle,color=blue];
      \node (X) at (2, 2) {$X$};
      \node (R) at (4, 2) {$R$};
      \node (Y) at (6, 2) {$Y$};
      \node (S) at (8, 2) {$S$};
      \node (Z) at (10, 2) {$Z$};
      \node (U1) at (4, 4) {$U_1$};
      \node (U2) at (8, 0) {$U_2$};
     \draw[-{Latex[length=3mm, width=2mm]}, color=blue] (X) to (R);
     \draw[-{Latex[length=3mm, width=2mm]}, color=blue] (R) to (Y);
     \draw[-{Latex[length=3mm, width=2mm]}, color=blue] (S) to (Y);
     \draw[-{Latex[length=3mm, width=2mm]}, color=blue] (Z) to (S);
     \draw[-{Latex[length=3mm, width=2mm]}, color=blue, dashed, bend right] (U1) to (X);
     \draw[-{Latex[length=3mm, width=2mm]}, color=blue, dashed, bend left] (U1) to (S);
     \draw[-{Latex[length=3mm, width=2mm]}, color=blue, dashed, bend left] (U2) to (R);
     \draw[-{Latex[length=3mm, width=2mm]}, color=blue, dashed, bend right] (U2) to (Z);
  \end{tikzpicture}
  \caption{The graph of an IL-CBN.}\label{fig:my_fig}
  \end{figure}

Look at the IL-CBN indicated by the graph in Figure \ref{fig:my_fig}.
We assume that $X, R, Y, S$, and $Z$ are observed, whereas $U_1$ and $U_2$ are latent.
Furthermore, we assume there are no input variables.
Do the following exercises, by using the chain-rule from above, the do-calculus, and your general understanding:
\begin{enumerate}
  \item Show that $P(Y |  \doit(X), \doit(Z)) = P(Y | R, S, \doit(X), \doit(Z)) \circ P(R, S | \doit(X), \doit(Z))$. 
  \item Show that $P(Y | R, S, \doit(X), \doit(Z)) = P(Y | R, S)$. 
  \item Show that $P(R , S| \doit(X), \doit(Z)) = P(R | S, \doit(X), \doit(Z)) \otimes P(S | \doit(X), \doit(Z))$.
  \item Show that $P(R | S, \doit(X), \doit(Z)) = P(R | X)$.
  \item Show that $P(S | \doit(X), \doit(Z)) =  P(S | Z)$.
  \item Argue that the causal effect $P(Y | \doit(X), \doit(Z))$ is identifiable by giving a formula in terms of observable probabilities.
\end{enumerate}

\section{Do-Calculus II}

Look at the IL-CBN indicated by the graph in Figure \ref{fig:my_second_fig}.
We assume that $X_1, X_2, X_3, X_4$, and $Y$ are observed, whereas $U_1$, $U_2$, and $U_3$ are latent.
Furthermore, we assume there are no input variables.
Do the following exercises, by using the chain-rule from above, the do-calculus, and your general understanding:

\begin{figure}[ht]
  \centering
  \begin{tikzpicture}[scale=1, transform shape]
      \tikzstyle{every node} = [draw,shape=circle,color=blue];
      \node (X1) at (2, 2) {$X_1$};
      \node (X2) at (4, 2) {$X_2$};
      \node (X3) at (6, 2) {$X_3$};
      \node (X4) at (8, 2) {$X_4$};
      \node (Y) at (10, 2) {$Y$};
      \node (U1) at (4, 0) {$U_1$};
      \node (U2) at (6, 4) {$U_2$};
      \node (U3) at (8, 0) {$U_3$};
     \draw[-{Latex[length=3mm, width=2mm]}, color=blue] (X1) to (X2);
     \draw[-{Latex[length=3mm, width=2mm]}, color=blue] (X2) to (X3);
     \draw[-{Latex[length=3mm, width=2mm]}, color=blue] (X3) to (X4);
     \draw[-{Latex[length=3mm, width=2mm]}, color=blue] (X4) to (Y);
     \draw[-{Latex[length=3mm, width=2mm]}, color=blue, dashed, bend left] (U1) to (X1);
     \draw[-{Latex[length=3mm, width=2mm]}, color=blue, dashed, bend right] (U1) to (X3);
     \draw[-{Latex[length=3mm, width=2mm]}, color=blue, dashed, bend right] (U2) to (X2);
     \draw[-{Latex[length=3mm, width=2mm]}, color=blue, dashed, bend left] (U2) to (X4);
     \draw[-{Latex[length=3mm, width=2mm]}, color=blue, dashed, bend left] (U3) to (X3);
     \draw[-{Latex[length=3mm, width=2mm]}, color=blue, dashed, bend right] (U3) to (Y);
  \end{tikzpicture}
  \caption{The graph of an IL-CBN}\label{fig:my_second_fig}
  \end{figure}

\begin{enumerate}
  \item Show that $P(Y | \doit(X_1, X_2, X_3, X_4)) = P(Y | \doit(X_4))$. 
  \item Show that $P(Y | \doit(X_4)) = P(Y | \doit(X_4), X_1, X_2, X_3) \circ P(X_1, X_2, X_3 | \doit(X_4))$.  
  \item Show that $P(Y | \doit(X_4), X_1, X_2, X_3) = P(Y | X_4, X_1, X_2, X_3)$. 
  \item Show that $P(X_1, X_2, X_3 | \doit(X_4)) = P(X_1, X_2, X_3)$.
  \item Argue that the causal effect $P(Y | \doit(X_1, X_2, X_3, X_4))$ is identifiable by giving a formula in terms of observationsl probabilities
  \end{enumerate}


\section{Structural Causal Models and Faithfulness}

Consider the following structural equations:
\begin{align*}
  X_1 & = E_1 + X_3, \\
  X_2 & = E_2 - X_1 + X_3 + X_4, \\ 
  X_3 & = E_3,  \\
  X_4 & = E_4 + X_1 - X_3. 
\end{align*}
Thereby, all variables are real valued.
The $E_i$ are exogenous random variables, i.i.d. standard univariate normal distributed, and the $X_i$ are endogenous variables.
There are no exogenous input variables.
Do the following:
\begin{enumerate}
  \item  {\color{magenta}Formally model the structural equations with an SCM $\mathcal{M}$.
    To simplify the notation, use the same symbols for the labels of the variables as for the variables themselves.}
\textbf{Did we want to remove part 1? I'm not sure anymore.}
  \item  Find a solution function $g$ for an SCM $\mathcal{M}$ representing these equations by finding an expression of the $X_i$ purely in terms of the $E_i$, such that the $X_i$ satisfy the structural equations. 
    Is it unique? If so, why? If not, find a second solution function. 
  \item  Draw the observational causal graph $G(\mathcal{M})^{\{X_1, X_2, X_3, X_4\}}$.
    {\color{magenta}
  \item  List all conditional independencies 
    \begin{equation*}
      X_i \Indep_{P(\boldsymbol{X}, \boldsymbol{E})} X_j | X_k
    \end{equation*} 
    for pairwise different $i, j, k$, where $\boldsymbol{X} = (X_1, X_2, X_3, X_4)$ and $\boldsymbol{E} = (E_1, E_2, E_3, E_4)$. 
    Argue that the conditional independencies which you find are all conditional independencies involving only single variables $X_i$.
  \item  List all $d$-separations 
    \begin{equation*}
      X_i \Perp^d_{G(\mathcal{M})} X_j | X_k
    \end{equation*}
    for pairwise different $i, j, k$, and argue that these are all such $d$-separations.
  \item  Argue that $\mathcal{M}$ is not faithful.
  \item Look at all the coefficients in front of any of the $X_i$ in any of the structural equations (all are either $+1$ or $-1$).
    Is it possible to change only one single coefficient to a different real number to get a faithful model? 
    If so, make this change. 
  If not, explain why this is not possible.} \textbf{Below: New version of 4, 5, 6, 7 which should be more solvable.}
\item Find three different independencies $X_i \Indep_{P(\boldsymbol{X}, \boldsymbol{E})} X_j$ with $\boldsymbol{X} = (X_1, X_2, X_3, X_4)$, $\boldsymbol{E} = (E_1, E_2, E_3, E_4)$ and $i < j$.
\item Is there a $d$-separation of the form  $X_i \Perp^d_{G(\mathcal{M})} X_j$ with $i, j \in \{1, 2, 3, 4\}$?
  If so, find it, of not, argue why not.
\item Is $\mathcal{M}$ faithful? Explain.
\item Look at all the coefficients in front of any of the $X_i$ in any of the structural equations (all are either $+1$ or $-1$).
  Is it possible to change only one single coefficient to a different real number to get a model such that all $X_i$ are \emph{dependent} on all $X_j$ for any $i, j \in \{1, 2, 3, 4\}$?
    If so, make this change. 
    If not, explain why this is not possible.
\end{enumerate}

\begin{gradedsolution}
  Part 11: If I'm not mistaken, then changing the coefficient in front of $X_1$ in the equation of $X_4$ to any number $\alpha \neq 1$ leads to all dependencies being true:
  \begin{equation*}
         X_4 = E_4 + \alpha X_1 - X_3
  \end{equation*}
\end{gradedsolution}


\section{$d$-separation cannot be formulated with paths instead of walks}

Construct a DAG $\mathcal{G}$ with four nodes $V = \{A, B, C, D\}$ such that the following two properties hold:
\begin{enumerate}
  \item  $\{A\} \nPerp_G^d \{B\} | \{C\}$. 
  \item  All \emph{paths}\footnote{Remember that, by definition, paths do not contain any node twice.} from $A$ to $B$ are $d$-blocked by $C$. 
\end{enumerate}

\begin{gradedsolution}

 \begin{figure}[ht]
  \centering
  \begin{tikzpicture}[scale=1, transform shape]
      \tikzstyle{every node} = [draw,shape=circle,color=blue];
      \node (B) at (2, 2) {$B$};
      \node (D) at (4, 2) {$D$};
      \node (A) at (6, 2) {$A$};
      \node (C) at (4, 0) {$C$};
     \draw[-{Latex[length=3mm, width=2mm]}, color=blue] (B) to (D);
     \draw[-{Latex[length=3mm, width=2mm]}, color=blue] (A) to (D);
     \draw[-{Latex[length=3mm, width=2mm]}, color=blue] (D) to (C);
  \end{tikzpicture}
  \caption{The caption of a possible solution}\label{fig:my_third_fig}
  \end{figure}

\end{gradedsolution}

\end{document}
